\section*{Bab \ref{ch:b1:logic} --- Logika, Himpunan dan Pemetaan}

\begin{solution}{exo:b1:logic:1}
Negasi, dengan mendorong $\lnot$ melewati setiap kuantor
(\cref{prop:b1:logic:negquant}) dan memakai $\lnot(P \implies Q) \iff
P \land \lnot Q$:
\begin{enumerate}
  \item $\exists x \in \R,\ \forall y \in \R,\ x + y \leq 0$;
  \item $\forall x \in \R,\ \exists y \in \R,\ xy \neq 0$;
  \item $\exists \varepsilon > 0,\ \forall \delta > 0,\ \exists x \in
        \R,\ \abs{x} \leq \delta \text{ dan } \abs{f(x)} >
        \varepsilon$.
\end{enumerate}
\omterm{def:b1:logic:statement}{Pernyataan} (1) benar: diberikan $x$, ambil $y = -x + 1$; maka $x + y = 1 >
0$. \omterm{def:b1:logic:statement}{Pernyataan} (2) benar: $x = 0$ memenuhi $xy = 0$ untuk setiap $y$.
\end{solution}

\begin{solution}{exo:b1:logic:2}
Tabel kebenaran, dengan menulis B/S untuk keempat kasus $(P, Q)$:
\begin{center}
\begin{tabular}{cc|c|c|c|c}
$P$ & $Q$ & $P \implies Q$ & $\lnot(P \implies Q)$ & $\lnot Q$ &
$P \land \lnot Q$ \\
\hline
B & B & B & S & S & S \\
B & S & S & B & B & B \\
S & B & B & S & S & S \\
S & S & B & S & B & S \\
\end{tabular}
\end{center}
Kolom $4$ dan $6$ berimpit, dan itu membuktikan ekuivalensinya. Negasi
dari ``jika sebuah fungsi dapat diturunkan maka ia kontinu'' karena itu
berbunyi: ``ada fungsi yang dapat diturunkan dan tidak kontinu''
(sebuah \omterm{def:b1:logic:statement}{pernyataan} yang salah, kebetulan saja: implikasi aslinya benar,
lihat \cref{ch:b1:derivative}).
\end{solution}

\begin{solution}{exo:b1:logic:3}
\emph{Kontraposisi.} Andaikan $x < 1$. Maka $x^3 < 1$ (fungsi pangkat
tiga naik) dan $x < 1$, sehingga $x^3 + x < 2$. Ini membuktikan
kontraposisinya, jadi juga \omterm{def:b1:logic:statement}{pernyataan} itu sendiri.

\emph{Kontradiksi.} Andaikan $a > 0$ bilangan real positif tegas yang
terkecil. Maka $a/2$ positif tegas dan $a/2 < a$ (karena $a > 0$),
bertentangan dengan keminimalannya. Jadi $a$ semacam itu tidak ada.
\end{solution}

\begin{solution}{exo:b1:logic:4}
\begin{enumerate}
  \item Basis $n = 0$: $2^0 = 1 = 2^1 - 1$. Langkah: dengan
        mengandaikan kesamaan itu berlaku untuk $n$,
        \[
        \sum_{k=0}^{n+1} 2^k = (2^{n+1} - 1) + 2^{n+1}
        = 2 \cdot 2^{n+1} - 1 = 2^{n+2} - 1 .
        \]
  \item Basis $n = 0$: $4^0 + 5 = 6 = 3 \times 2$. Langkah: jika
        $4^n + 5 = 3m$, maka
        \[
        4^{n+1} + 5 = 4(4^n + 5) - 15 = 3(4m - 5),
        \]
        yang habis dibagi $3$.
\end{enumerate}
\end{solution}

\begin{solution}{exo:b1:logic:5}
Langkah induksinya diam-diam mengandaikan bahwa kedua kelompok
(``$n$ pertama'' dan ``$n$ terakhir'') saling tumpang-tindih, sehingga
pensil bersamanya membawa warna dari kelompok yang satu ke kelompok
yang lain. Untuk $n + 1 = 2$ kedua kelompok itu adalah
$\{$pensil pertama$\}$ dan $\{$pensil kedua$\}$: keduanya saling lepas,
dan argumennya runtuh. Jadi $P(1) \implies P(2)$ tidak pernah
dibuktikan, dan induksinya ambruk --- walaupun $P(n) \implies P(n+1)$
sah untuk setiap $n \geq 2$.
\end{solution}

\begin{solution}{exo:b1:logic:6}
\begin{enumerate}
  \item $x \in A \setminus B \iff x \in A \land x \notin B \iff x \in
        A \land x \in \overline{B} \iff x \in A \cap \overline{B}$.
  \item Dengan memakai (1) dan sifat distributif
        (\cref{prop:b1:logic:setalgebra}):
        $(A \cup B) \cap \overline{C} = (A \cap \overline{C}) \cup
        (B \cap \overline{C})$.
  \item Andaikan $A \subseteq B$. Maka $A \cup B \subseteq B$ (kedua
        potongannya terletak di $B$) dan $B \subseteq A \cup B$ selalu
        berlaku, sehingga $A \cup B = B$. Andaikan $A \cup B = B$: maka
        $A \cap B
        \subseteq A$ selalu berlaku, dan $A \subseteq A \cup B = B$
        memberikan $A \subseteq A \cap B$, sehingga $A \cap B = A$.
        Andaikan $A \cap B =
        A$: maka $A = A \cap B \subseteq B$. Jadi ketiga syarat itu
        ekuivalen (kita membuktikan satu daur implikasi).
\end{enumerate}
\end{solution}

\begin{solution}{exo:b1:logic:7}
\begin{enumerate}
  \item $f$ \omterm{def:b1:logic:inj}{injektif} ($n + 1 = m + 1 \implies n = m$) tetapi tidak
        \omterm{def:b1:logic:inj}{surjektif}: $0$ tak mempunyai \omterm{def:b1:logic:map}{prapeta} di $\N$.
  \item $g$ \omterm{def:b1:logic:inj}{bijektif}: $n \mapsto n - 1$ adalah invers dua sisi pada
        $\Z$.
  \item $h$ \omterm{def:b1:logic:inj}{injektif}: $\frac{x+1}{x-1} = \frac{x'+1}{x'-1}$ memberikan
        $(x+1)(x'-1) = (x'+1)(x-1)$, yaitu $xx' - x + x' - 1 = xx' -
        x' + x - 1$, sehingga $2x' = 2x$. Ia tidak \omterm{def:b1:logic:inj}{surjektif} pada $\R$:
        menyelesaikan $y = \frac{x+1}{x-1}$ memberikan $x(y - 1) = y + 1$,
        yang tak punya penyelesaian bila $y = 1$ (persamaannya menjadi
        $0 = 2$). Dengan daerah kawan $\R \setminus \{1\}$, perhitungan
        yang sama memberikan \omterm{def:b1:logic:map}{prapeta} tunggal $x = \frac{y+1}{y-1}$, sehingga
        $h \colon \R \setminus \{1\} \to \R \setminus \{1\}$ bersifat
        \omterm{def:b1:logic:inj}{bijektif} dan $h^{-1}(y) = \frac{y+1}{y-1} = h(y)$: $h$ adalah
        inversnya sendiri.
\end{enumerate}
\end{solution}

\begin{solution}{exo:b1:logic:8}
\begin{enumerate}
  \item $x \in f^{-1}(B \cap B') \iff f(x) \in B \cap B' \iff f(x) \in
        B \land f(x) \in B' \iff x \in f^{-1}(B) \cap f^{-1}(B')$.
        Untuk peta: $y \in f(A \cup A')$ berlaku persis bila $y = f(x)$ untuk suatu $x$
        di $A$ atau di $A'$, yakni bila $y \in f(A)$ atau $y \in f(A')$.
  \item Jika $y \in f(A \cap A')$, maka $y = f(x)$ dengan $x \in A$ dan
        $x \in A'$, sehingga $y \in f(A)$ dan $y \in f(A')$. Kesejatiannya:
        ambil $f \colon \R \to \R$, $x \mapsto x^2$, $A = \{-1\}$,
        $A' = \{1\}$: maka $f(A \cap A') = f(\emptyset) = \emptyset$
        tetapi $f(A) \cap f(A') = \{1\}$.
  \item ($\Leftarrow$) Dengan $A = \{x\}$, $A' = \{x'\}$ untuk $x \neq
        x'$: jika $f(x) = f(x')$, maka $f(A) \cap f(A') = \{f(x)\}$
        sedangkan $f(A \cap A') = \emptyset$, bertentangan dengan
        kesamaan yang diandaikan; jadi $f$ \omterm{def:b1:logic:inj}{injektif}. ($\Rightarrow$) Misalkan $f$
        \omterm{def:b1:logic:inj}{injektif} dan $y \in f(A) \cap f(A')$: $y = f(x) = f(x')$ dengan
        $x \in A$, $x' \in A'$; keinjektifan memberikan $x = x' \in A \cap
        A'$, sehingga $y \in f(A \cap A')$. Bersama (2), kesamaannya berlaku.
\end{enumerate}
\end{solution}

\begin{solution}{exo:b1:logic:9}
$g \circ f = \mathrm{id}_E$ bersifat \omterm{def:b1:logic:inj}{injektif} sekaligus \omterm{def:b1:logic:inj}{surjektif}, jadi menurut
\cref{prop:b1:logic:comp}~(2), $f$ \omterm{def:b1:logic:inj}{injektif} dan $g$ \omterm{def:b1:logic:inj}{surjektif}.
Contoh: $E = \N$, $F = \Z$, $f$ \omterm{def:b1:logic:map}{pemetaan} inklusi $n \mapsto n$, dan
$g \colon \Z \to \N$, $g(n) = n$ untuk $n \geq 0$ dan $g(n) = 0$ untuk
$n < 0$. Maka $g(f(n)) = n$ untuk setiap $n \in \N$, tetapi $f$ tidak
\omterm{def:b1:logic:inj}{surjektif} dan $g$ tidak \omterm{def:b1:logic:inj}{injektif}.
\end{solution}

\begin{solution}{exo:b1:logic:10}
$x^2 - y^2 = x - y \iff (x - y)(x + y) = x - y \iff (x - y)(x + y - 1)
= 0 \iff y = x$ atau $y = 1 - x$. \emph{Refleksif:} $y = x$ memenuhi.
\emph{Simetris:} syarat ``$y = x$ atau $y = 1 - x$'' bersifat
simetris dalam $x$ dan $y$ (jika $y = 1 - x$ maka $x = 1 - y$).
\emph{Transitif:} andaikan $x \mathbin{\mathcal{R}} y$ dan $y
\mathbin{\mathcal{R}} z$; dengan menelusuri keempat kasusnya, $z$ selalu sama dengan $x$
atau $1 - x$ (misalnya $y = 1 - x$ dan $z = 1 - y$ memberikan $z = x$).
Jadi $\mathcal{R}$ adalah \omterm{def:b1:logic:equiv}{relasi ekuivalensi} dan $\mathrm{cl}(x) =
\{x,\, 1 - x\}$. Kelas ini beranggotakan satu unsur tepat ketika $x = 1 - x$,
yaitu untuk $x = \frac12$.
\end{solution}

\begin{solution}{exo:b1:logic:11}
Misalkan $f \colon E \to \mathcal{P}(E)$ sebuah \omterm{def:b1:logic:map}{pemetaan} sembarang lalu tulis
$D = \{x \in E : x \notin f(x)\} \in \mathcal{P}(E)$. Andaikan $D =
f(a)$ untuk suatu $a \in E$. Jika $a \in D$, maka menurut definisi $D$,
$a \notin f(a) = D$: kontradiksi. Jika $a \notin D$, maka $a \notin
f(a)$, sehingga menurut definisi $D$, $a \in D$: kontradiksi. Jadi $D$
tidak berada pada peta $f$, dan $f$ tidak \omterm{def:b1:logic:inj}{surjektif}. (Khususnya tak ada
\omterm{def:b1:logic:sets}{himpunan} yang berbijeksi dengan \omterm{def:b1:logic:sets}{himpunan kuasanya}: ada ``lebih banyak''
\omterm{def:b1:logic:sets}{himpunan} bagian $\N$ daripada bilangan bulat.)
\end{solution}

\begin{solution}{exo:b1:logic:12}
\begin{enumerate}
  \item ($\Rightarrow$) Misalkan $f$ \omterm{def:b1:logic:inj}{surjektif} dan $\Phi(B) =
        \Phi(B')$. Untuk $y \in B$, pilih $x$ dengan $f(x) = y$; maka $x
        \in f^{-1}(B) = f^{-1}(B')$, sehingga $y = f(x) \in B'$. Jadi $B
        \subseteq B'$, dan secara simetris $B' \subseteq B$: $\Phi$
        \omterm{def:b1:logic:inj}{injektif}. ($\Leftarrow$) Jika $f$ tidak \omterm{def:b1:logic:inj}{surjektif}, pilih $y_0
        \in F$ di luar petanya; maka $f^{-1}(\{y_0\}) = \emptyset =
        f^{-1}(\emptyset)$ padahal $\{y_0\} \neq \emptyset$, sehingga $\Phi$
        tidak \omterm{def:b1:logic:inj}{injektif}.
  \item ($\Rightarrow$) Misalkan $f$ \omterm{def:b1:logic:inj}{injektif} dan $A \subseteq E$. Tulis
        $B = f(A)$; maka $f^{-1}(B) = \{x : f(x) \in f(A)\}$, dan
        keinjektifan memberikan $f(x) \in f(A) \iff x \in A$, sehingga $\Phi(B) =
        A$: jadi $\Phi$ \omterm{def:b1:logic:inj}{surjektif}. ($\Leftarrow$) Jika $f$ tidak
        \omterm{def:b1:logic:inj}{injektif}, ambil $x \neq x'$ dengan $f(x) = f(x')$. Setiap
        \omterm{def:b1:logic:sets}{himpunan} \omterm{def:b1:logic:map}{prapeta} $f^{-1}(B)$ memuat $x$ jika dan hanya jika ia
        memuat $x'$; jadi $\{x\}$ tidak berbentuk $\Phi(B)$, dan
        $\Phi$ tidak \omterm{def:b1:logic:inj}{surjektif}.
\end{enumerate}
\end{solution}

\begin{solution}{pb:b1:logic:1}
\textbf{1.} \emph{Refleksif:} $\mathrm{id}_E$ adalah bijeksi $E$
pada dirinya sendiri. \emph{Simetris:} jika $f \colon E \to F$ \omterm{def:b1:logic:inj}{bijektif},
\cref{thm:b1:logic:inverse} menyediakan $f^{-1} \colon F \to E$, yang juga
\omterm{def:b1:logic:inj}{bijektif}. \emph{Transitif:} jika $f \colon E \to F$ dan $g \colon F
\to G$ bijeksi, \cref{prop:b1:logic:comp}~(1) mengatakan $g \circ f
\colon E \to G$ bijeksi. (Ini hanya ``seperti'' \omterm{def:b1:logic:equiv}{relasi
ekuivalensi}: kumpulan semua \omterm{def:b1:logic:sets}{himpunan} terbukanlah \omterm{def:b1:logic:sets}{himpunan}, karena
paradoks yang disinggung \cref{exo:b1:logic:11}; ketiga sifat itulah
yang penting.)

\textbf{2.} Jika $f \colon E \to F$ dan $g \colon F \to G$
\omterm{def:b1:logic:inj}{injektif}, maka $g \circ f$ \omterm{def:b1:logic:inj}{injektif} menurut \cref{prop:b1:logic:comp}~(1):
$E \preceq G$. Untuk butir kedua, batasi daerah kawan $f$ pada petanya:
\omterm{def:b1:logic:map}{pemetaan} $\tilde f \colon E \to f(E)$, $x \mapsto f(x)$, \omterm{def:b1:logic:inj}{surjektif} menurut
konstruksi $f(E)$ dan \omterm{def:b1:logic:inj}{injektif} karena $f$ \omterm{def:b1:logic:inj}{injektif}, jadi \omterm{def:b1:logic:inj}{bijektif}:
$E \approx f(E)$.

\textbf{3.} ($\Rightarrow$) Misalkan $f \colon E \to F$ \omterm{def:b1:logic:inj}{injektif} dan
tetapkan $a \in E$ ($E \neq \emptyset$). Definisikan $s \colon F \to E$ oleh:
$s(y)$ adalah satu-satunya $x$ dengan $f(x) = y$ bila $y \in f(E)$
(ketunggalannya menurut keinjektifan), dan $s(y) = a$ bila tidak. Untuk setiap $x
\in E$, $s(f(x)) = x$, jadi setiap $x$ tercapai: $s$ \omterm{def:b1:logic:inj}{surjektif}.
($\Leftarrow$) Misalkan $s \colon F \to E$ \omterm{def:b1:logic:inj}{surjektif}. Untuk masing-masing $x \in
E$ pilih satu $y_x \in F$ dengan $s(y_x) = x$, lalu tulis $u(x) = y_x$. Jika
$u(x) = u(x')$ maka $x = s(u(x)) = s(u(x')) = x'$: jadi $u \colon E \to F$
\omterm{def:b1:logic:inj}{injektif}.

\textbf{4.} $n \mapsto n + 1$ memetakan $\N$ ke dalam $\N^*$, bersifat \omterm{def:b1:logic:inj}{injektif}
($n + 1 = m + 1 \implies n = m$) dan \omterm{def:b1:logic:inj}{surjektif} (setiap $m \geq 1$ sama dengan
$(m - 1) + 1$ dengan $m - 1 \in \N$). Untuk $\sigma$: ia memetakan bilangan
genap $0, 2, 4, \dots$ ke $0, 1, 2, \dots$ dan bilangan ganjil $1, 3,
5, \dots$ ke $-1, -2, -3, \dots$ Keinjektifan: masukan genap mendarat di
$\N$ ($\sigma(n) = n/2 \geq 0$) dan masukan ganjil mendarat di
bilangan bulat negatif tegas ($\sigma(n) = -(n+1)/2 \leq -1$), jadi
tabrakan hanya mungkin terjadi di dalam satu kelas paritas, tempat $\sigma$
monoton tegas ($n/2 = m/2$ atau $(n+1)/2 = (m+1)/2$ memaksa $n =
m$). Kesurjektifan: $k \geq 0$ sama dengan $\sigma(2k)$; $k \leq -1$ sama dengan
$\sigma(-2k - 1)$ dengan $-2k - 1 \geq 1$ ganjil. Jadi $\N \approx \N^*$ dan
$\N \approx \Z$.

\textbf{5.} $C_0 = E \setminus g(F) \subseteq C$, jadi $x \notin C$
mengakibatkan $x \notin C_0$, yaitu $x \in g(F)$: ada $y \in F$ yang memenuhi
$g(y) = x$. Jika juga $g(y') = x$, keinjektifan $g$ memberikan $y' = y$.
Jadi klausa kedua pada definisi $h$ memilih satu unsur $g^{-1}(x)$ yang
tunggal dan terdefinisi dengan baik.

\textbf{6.} Peta langsung berkomutasi dengan gabungan
(\cref{exo:b1:logic:8}~(1), diterapkan pada $f$ lalu pada $g$):
\[
g\bigl(f(C)\bigr)
= g\Bigl(f\Bigl(\bigcup_{n \in \N} C_n\Bigr)\Bigr)
= \bigcup_{n \in \N} g\bigl(f(C_n)\bigr)
= \bigcup_{n \in \N} C_{n+1}
= \bigcup_{n \geq 1} C_n \subseteq C .
\]

\textbf{7.} Ambil $x \neq x'$ di $E$. Jika keduanya terletak di $C$, maka $h(x) =
f(x) \neq f(x') = h(x')$ menurut keinjektifan $f$. Jika keduanya tidak terletak di
$C$, maka $g(h(x)) = x \neq x' = g(h(x'))$, sehingga $h(x) \neq h(x')$. Jika
$x \in C$ dan $x' \notin C$ (kasus campuran, sampai penukaran nama):
andaikan $h(x) = h(x')$, yaitu $f(x) = g^{-1}(x')$. Dengan menerapkan $g$:
$x' = g(f(x)) \in g(f(C))$, dan pertanyaan 6 memberikan $x' \in C$ ---
kontradiksi. Jadi $h(x) \neq h(x')$ pada semua kasus: $h$ \omterm{def:b1:logic:inj}{injektif}.

\textbf{8.} Ambil $y \in F$. \emph{Kasus 1: $g(y) \notin C$.} Maka
$h(g(y)) = g^{-1}(g(y)) = y$: unsur $g(y)$ adalah sebuah \omterm{def:b1:logic:map}{prapeta}.
\emph{Kasus 2: $g(y) \in C$}, katakanlah $g(y) \in C_n$. Karena $g(y) \in
g(F)$, kita punya $g(y) \notin C_0 = E \setminus g(F)$, jadi $n \geq 1$
dan $g(y) \in C_n = g(f(C_{n-1}))$: ada $x \in C_{n-1}$ dengan
$g(y) = g(f(x))$. Keinjektifan $g$ memberikan $y = f(x)$, dan $x \in
C_{n-1} \subseteq C$, sehingga $h(x) = f(x) = y$. Pada kedua kasus $y$
tercapai: $h$ \omterm{def:b1:logic:inj}{surjektif}, jadi \omterm{def:b1:logic:inj}{bijektif}.

\textbf{9.} Jika $E \preceq F$ dan $F \preceq E$, pilih injeksi $f
\colon E \to F$ dan $g \colon F \to E$; pertanyaan 5--8 membangun sebuah
bijeksi $h \colon E \to F$, sehingga $E \approx F$. \omterm{def:b1:logic:statement}{Pernyataan} itu
tidak sepele karena kedua injeksi yang diberikan tak saling berkaitan --- tak satu pun
harus \omterm{def:b1:logic:inj}{surjektif}, dan tak ada rumus naif yang mencampur $f$ dan $g$ yang mendefinisikan
sebuah \omterm{def:b1:logic:map}{pemetaan}: seluruh isinya adalah \omterm{thm:b1:logic:partition}{partisi} $E$ atas daerah $C$
(tempat kita menyalin $f$) dan komplemennya (tempat kita menjalankan $g$
secara terbalik).

\textbf{10.} (a) \omterm{def:b1:logic:map}{Pemetaan} inklusi $\intoo01 \to \intcc01$ bersifat \omterm{def:b1:logic:inj}{injektif};
dan $x \mapsto \frac{x + 1}3$ memetakan $\intcc01$ secara \omterm{def:b1:logic:inj}{injektif} ke dalam
$\intcc{\frac13}{\frac23} \subseteq \intoo01$ (ia afin dengan
kemiringan tak nol). Menurut pertanyaan 9, $\intcc01 \approx \intoo01$ --- sebuah
bijeksi yang cukup menjengkelkan untuk dituliskan secara eksplisit.
(b) \emph{Keinjektifan.} Andaikan $2^p(2q + 1) = 2^{p'}(2q' + 1)$ dengan,
katakanlah, $p \leq p'$. Dengan membagi $2^p$: $2q + 1 = 2^{p' - p}(2q' + 1)$.
Jika $p' > p$ ruas kanan genap dan ruas kiri ganjil ---
mustahil; jadi $p = p'$, lalu $2q + 1 = 2q' + 1$ dan $q = q'$.
\emph{Kesurjektifan.} Kita tunjukkan dengan induksi kuat bahwa setiap bilangan bulat
$m \geq 1$ berbentuk $2^p(2q + 1)$. Untuk $m = 1$: $p = q = 0$.
Ambil $m \geq 1$ dan andaikan klaim itu berlaku untuk semua bilangan bulat di
$\intint1m$. Jika $m + 1$ ganjil, $m + 1 = 2q + 1$ dengan $p = 0$. Jika
$m + 1$ genap, $m + 1 = 2m'$ dengan $1 \leq m' \leq m$; menurut
hipotesis $m' = 2^p(2q + 1)$, sehingga $m + 1 = 2^{p+1}(2q + 1)$. Jadi
$\varphi(p, q) = 2^p(2q + 1) - 1$ mengenai setiap $n \in \N$, dan
$\varphi$ adalah bijeksi $\N \times \N \to \N$.

\textbf{11.} Karena $A$ tak hingga, $A \setminus \{\varphi(0), \dots,
\varphi(n - 1)\}$ tak pernah kosong, dan sifat unsur terkecil pada
$\N$ (yang dipakai untuk membuktikan \cref{thm:b1:logic:induction}) membuat
definisi rekursif itu sah. \emph{Naik tegas:}
$\varphi(n + 1)$ termasuk $A \setminus \{\varphi(0), \dots,
\varphi(n)\} \subseteq A \setminus \{\varphi(0), \dots, \varphi(n -
1)\}$, yang minimumnya adalah $\varphi(n)$; jadi $\varphi(n + 1) \geq
\varphi(n)$, dan kesamaan tersingkir, sehingga $\varphi(n+1) >
\varphi(n)$. \emph{$\varphi(n) \geq n$:} dengan induksi, $\varphi(0)
\geq 0$, dan $\varphi(n + 1) \geq \varphi(n) + 1 \geq n + 1$.
\emph{Keinjektifan} menyusul dari kemonotonan tegas.
\emph{Kesurjektifan pada $A$:} andaikan ada $a \in A$ yang tak pernah
tercapai. Karena $\varphi(a + 1) \geq a + 1 > a$, \omterm{def:b1:logic:sets}{himpunan} semua $n$ dengan
$\varphi(n) > a$ tak kosong; misalkan $n$ unsur terkecilnya. Untuk setiap
$k < n$ berlaku $\varphi(k) \leq a$, jadi $\varphi(k) < a$ ($a$ tidak
tercapai). Maka $a$ terletak di $A \setminus \{\varphi(0), \dots,
\varphi(n - 1)\}$ dan $a < \varphi(n)$, bertentangan dengan keminimalan
yang mendefinisikan $\varphi(n)$. Jadi $\varphi$ adalah bijeksi $\N \to A$, dan
$A$ terbilang.

\textbf{12.} Misalkan $E \preceq \N$ lewat injeksi $f$; maka $E \approx
f(E)$ (pertanyaan 2). Jika $f(E)$ hingga, $E$ hingga; jika $f(E)$
tak hingga, pertanyaan 11 memberikan $f(E) \approx \N$, sehingga $E \approx \N$ menurut
ketransitifan (pertanyaan 1). Sebaliknya \omterm{def:b1:logic:sets}{himpunan} hingga dan \omterm{pb:b1:logic:1}{himpunan terbilang}
jelas terinjeksi ke dalam $\N$. Jalan pintasnya: $E \preceq \N$ dan $\N
\preceq E$ memberikan $E \approx \N$ secara langsung menurut
Cantor--Schröder--Bernstein --- tanpa perlu argumen pencacahan.

\textbf{13.} Misalkan $f \colon E \to \N$ dan $g \colon F \to \N$ dua
injeksi. Maka $(x, y) \mapsto \varphi\bigl(f(x), g(y)\bigr)$ adalah
injeksi $E \times F \to \N$: jika petanya berimpit, keinjektifan
$\varphi$ (pertanyaan 10) memberikan $f(x) = f(x')$ dan $g(y) = g(y')$,
lalu $x = x'$, $y = y'$. Untuk $\Z \times \N^*$: kedua faktornya
terbilang (pertanyaan 4), jadi $\Z \times \N^* \preceq \N$; ia
tak hingga (ia memuat $\{0\} \times \N^*$), sehingga terbilang menurut
pertanyaan 12.

\textbf{14.} Setiap bilangan rasional $r$ mempunyai penyajian tunggal $r =
p/q$ dengan $p \in \Z$, $q \in \N^*$ dan pecahan itu paling sederhana
(ketunggalannya dibuktikan pada \cref{ch:b1:arith}; untuk $r = 0$ ambil
$0/1$). \omterm{def:b1:logic:map}{Pemetaan} $r \mapsto (p, q)$ lalu bersifat \omterm{def:b1:logic:inj}{injektif}: pasangan itu
menentukan $r = p/q$. Jadi $\Q \preceq \Z \times \N^* \preceq \N$
menurut pertanyaan 13. Karena $\N \subseteq \Q$ memberikan $\N \preceq \Q$,
pertanyaan 12 (atau langsung Cantor--Schröder--Bernstein) menunjukkan $\Q
\approx \N$: bilangan rasional bersifat terbilang.

\textbf{15.} Untuk masing-masing $n$ tetapkan injeksi $f_n \colon E_n \to \N$.
Untuk $x \in \bigcup_n E_n$, misalkan $n(x)$ indeks $n$ \emph{terkecil} yang memenuhi
$x \in E_n$, lalu tulis $u(x) = \varphi\bigl(n(x), f_{n(x)}(x)\bigr)
\in \N$. Jika $u(x) = u(x')$, keinjektifan $\varphi$ memberikan $n(x) =
n(x') = n$ dan $f_n(x) = f_n(x')$, sehingga $x = x'$ menurut keinjektifan
$f_n$. Jadi gabungan itu terinjeksi ke dalam $\N$: ia paling banyak terbilang.

\textbf{16.} Tulis $\Psi(F) = \sum_{i \in F} 2^i$ untuk $F \subseteq \N$
yang hingga ($\Psi(\emptyset) = 0$). Andaikan $F \neq F'$ dan misalkan $m$
unsur terbesar tempat keduanya berbeda, katakanlah $m \in F \setminus
F'$ (tukar namanya jika perlu). Unsur yang $> m$ termasuk keduanya atau
tak termasuk keduanya, jadi sumbangannya sama pada kedua jumlah itu;
dengan membandingkan sumbangan unsur yang $\leq m$:
\[
\sum_{i \in F,\, i \leq m} 2^i \geq 2^m
> 2^m - 1 = \sum_{k=0}^{m-1} 2^k
\geq \sum_{i \in F',\, i \leq m} 2^i ,
\]
memakai jumlah geometri \cref{exo:b1:logic:4}. Jadi $\Psi(F)
\neq \Psi(F')$: $\Psi$ \omterm{def:b1:logic:inj}{injektif} dan \omterm{def:b1:logic:sets}{himpunan} semua \omterm{def:b1:logic:sets}{himpunan} bagian hingga
dari $\N$ paling banyak terbilang; ia tak hingga (ia memuat semua
\omterm{def:b1:logic:sets}{himpunan} beranggota tunggal), jadi terbilang.

\textbf{17.} Kirim $A \subseteq \N$ ke fungsi indikatornya $\mathbf 1_A
\colon \N \to \{0,1\}$, $\mathbf 1_A(n) = 1$ bila $n \in A$ dan $0$
bila tidak; kirim $u \in \{0,1\}^{\N}$ ke $A_u = \{n \in \N : u(n) =
1\}$. Kedua \omterm{def:b1:logic:map}{pemetaan} itu saling invers: $A_{\mathbf 1_A} = A$ dan
$\mathbf 1_{A_u} = u$ (periksa nilainya pada setiap $n$). Menurut
\cref{thm:b1:logic:inverse}, masing-masing adalah bijeksi: $\mathcal{P}(\N)
\approx \{0,1\}^{\N}$.

\textbf{18.} Untuk setiap $n$, $d(n) = 1 - \Phi(n)(n) \neq \Phi(n)(n)$,
jadi barisan $d$ dan $\Phi(n)$ berbeda pada indeks $n$: $d \neq
\Phi(n)$. Jadi tak ada $\Phi$ yang \omterm{def:b1:logic:inj}{surjektif}, dan menurut pertanyaan 3 tak ada pula
injeksi $\{0,1\}^{\N} \to \N$: jadi $\{0,1\}^{\N}$ bukan \omterm{def:b1:logic:sets}{himpunan} yang
paling banyak terbilang. Lewat kamus pertanyaan 17, sebuah \omterm{def:b1:logic:map}{pemetaan} $\Phi
\colon \N \to \{0,1\}^{\N}$ adalah \omterm{def:b1:logic:map}{pemetaan} $f \colon \N \to
\mathcal{P}(\N)$, dan $d$ berpadanan dengan \omterm{def:b1:logic:sets}{himpunan} $D = \{n : n
\notin f(n)\}$ (memang $d(n) = 1 \iff \Phi(n)(n) = 0 \iff n \notin
f(n)$): argumen diagonal itu \emph{adalah} bukti Cantor untuk
\cref{exo:b1:logic:11} dengan $E = \N$.

\textbf{19.} Tulis $x_n = 0.d_1(n)\,d_2(n)\,d_3(n)\dots$ dalam bentuk
sejati lalu definisikan $\delta_n = 5$ bila $d_n(n) \neq 5$, $\delta_n = 6$
bila $d_n(n) = 5$, lalu $x = 0.\delta_1\delta_2\delta_3\dots$ Ekspansi
ini hanya memakai angka $5$ dan $6$, jadi ia tidak berakhir dengan
angka $9$ semua: ia ekspansi sejati dari suatu bilangan real $x \in \intco01$.
Untuk setiap $n$, angka ke-$n$ dari $x$ dan $x_n$ berbeda
($\delta_n \neq d_n(n)$ menurut konstruksinya); karena ekspansi sejati
bersifat tunggal, $x \neq x_n$. Jadi tak ada barisan yang menghabiskan $\intco01$: menurut
pertanyaan 3 lagi, $\intco01$ bukan \omterm{def:b1:logic:sets}{himpunan} yang paling banyak terbilang.

\textbf{20.} $\intco01 \subseteq \R$, jadi injeksi $\R \to \N$
akan terbatasi menjadi injeksi pada $\intco01$, bertentangan dengan pertanyaan 19:
$\R$ tak terbilang. Jika $\R \setminus \Q$ paling banyak terbilang,
maka $\R = \Q \cup (\R \setminus \Q)$ akan menjadi gabungan dua \omterm{def:b1:logic:sets}{himpunan} yang paling
banyak terbilang, jadi paling banyak terbilang menurut pertanyaan 15 (ambil
$E_0 = \Q$, $E_n = \R \setminus \Q$ untuk $n \geq 1$) ---
kontradiksi. Jadi bilangan irasional tak terbilang. Persisnya:
di dalam $\R$, bilangan rasional membentuk \omterm{pb:b1:logic:1}{himpunan terbilang} sedangkan
komplemennya tak terbilang; tak ada bijeksi yang dapat memasangkan $\R
\setminus \Q$ dengan $\Q$ --- bilangan irasional secara tegas ``lebih banyak''
daripada bilangan rasional, meskipun keduanya tak hingga dan keduanya padat.

\textbf{21.} $p/q$ (dengan $q \neq 0$) adalah akar $qX - p$, sebuah
polinomial tak nol berkoefisien bulat. $\sqrt 2$ adalah akar
$X^2 - 2$. Untuk $x = \sqrt 2 + \sqrt 3$: $x^2 = 5 + 2\sqrt 6$, sehingga
$x^2 - 5 = 2\sqrt 6$ dan $(x^2 - 5)^2 = 24$, yaitu
\[
x^4 - 10x^2 + 1 = 0 :
\]
$\sqrt 2 + \sqrt 3$ adalah akar $X^4 - 10X^2 + 1$.

\textbf{22.} Petakan $P = a_0 + a_1X + \dots + a_nX^n$ (berderajat $\leq n$,
berkoefisien bulat) ke $(a_0, \dots, a_n) \in \Z^{n+1}$: \omterm{def:b1:logic:map}{pemetaan} ini
\omterm{def:b1:logic:inj}{injektif}, karena polinomial ditentukan oleh koefisiennya. Dengan
induksi pada $n$: $\Z^1 = \Z$ terbilang (pertanyaan 4), dan
$\Z^{n+2} \approx \Z^{n+1} \times \Z$ paling banyak terbilang menurut
pertanyaan 13. Jadi setiap \omterm{def:b1:logic:sets}{himpunan} polinomial bulat yang derajatnya terbatas
bersifat paling banyak terbilang; ia tak hingga (ia memuat semua konstanta), jadi
terbilang menurut pertanyaan 12.

\textbf{23.} \omterm{def:b1:logic:sets}{Himpunan} semua polinomial bulat adalah $\bigcup_{n \in
\N} \{P : \deg P \leq n,\ P \text{ berkoefisien bulat}\}$, sebuah
gabungan terbilang dari \omterm{pb:b1:logic:1}{himpunan terbilang}: paling banyak terbilang menurut
pertanyaan 15, tak hingga, jadi terbilang.

\textbf{24.} Untuk setiap polinomial bulat tak nol $P$, \omterm{def:b1:logic:sets}{himpunan} akarnya
$R_P = \{x \in \R : P(x) = 0\}$ bersifat hingga (paling banyak $\deg P$
unsur, diterima tanpa bukti). Menurut pertanyaan 23 polinomial bulat tak nol
dapat dicacah $P_0, P_1, P_2, \dots$; maka $\mathcal{A} =
\bigcup_{n \in \N} R_{P_n}$ adalah gabungan terbilang dari \omterm{def:b1:logic:sets}{himpunan} hingga (jadi
paling banyak terbilang): paling banyak terbilang menurut pertanyaan 15. Ia
memuat $\Q$ (pertanyaan 21), jadi ia tak hingga: $\mathcal{A}$
terbilang.

\textbf{25.} Jika $\R \setminus \mathcal{A}$ paling banyak terbilang,
$\R = \mathcal{A} \cup (\R \setminus \mathcal{A})$ akan paling banyak
terbilang (pertanyaan 15), bertentangan dengan pertanyaan 20. Jadi
\omterm{pb:b1:logic:1}{bilangan transenden} itu ada dan bahkan membentuk \omterm{def:b1:logic:sets}{himpunan} tak terbilang,
sedangkan \omterm{pb:b1:logic:1}{bilangan aljabar} --- yang mencakup setiap bilangan yang dibangun
dari bilangan bulat lewat penarikan akar --- hanya membentuk kerangka
terbilang di dalam $\R$. Ringkasan arsitekturnya: pertanyaan 1--3 menyiapkan
bahasa perbandingannya; Cantor--Schröder--Bernstein (pertanyaan
5--9) memungkinkan kita membuktikan ekuipotensi lewat dua injeksi mudah, bukan
lewat satu bijeksi cerdik, dan dipakai untuk $\intcc01 \approx \intoo01$,
untuk $\Q$ dan di sepanjang Bagian V; bijeksi pemasangan (pertanyaan
10) menggerakkan hasil kali dan gabungan terbilang (pertanyaan 13, 15),
yang pada gilirannya menggerakkan $\Q$, polinomial bulat dan
$\mathcal{A}$; argumen diagonal (pertanyaan 18--19) menyediakan
satu-satunya ketaksamaan tegas $\N \prec \R$ yang membuat seluruh kisah ini
tidak sepele. Kesimpulan Cantor mencengangkan secara filosofis: buktinya
sama sekali tidak menunjukkan \omterm{pb:b1:logic:1}{bilangan transenden} mana pun, namun menunjukkan bahwa dalam
pengertian ekuipotensi \emph{hampir setiap} bilangan real bersifat
transenden. Menyebut satu \omterm{pb:b1:logic:1}{bilangan transenden} tertentu --- $\pi$ atau
$\eu$ --- menuntut matematika yang sama sekali lain dan puluhan tahun
kerja tambahan.

\end{solution}
